\documentclass[12pt]{article}
\usepackage[utf8]{inputenc}
\usepackage[T1]{fontenc}
\usepackage{hyperref}
\title{Wik-Me'anha Language Resources for K-12 Curriculum Development}
\author{Compiled by Claude Sonnet 5.5 (claude-sonnet-5-5) for Andr\'as Kornai}
\date{}
\begin{document}
\maketitle

\section{Introduction}
Wik-Me'anha (ISO 639-3: wih; Glottocode: wikm1246), also written Wik-Me'nh or Wik Me'anh, is a Paman (Pama-Nyungan) language of Cape York Peninsula, Queensland. The only speaker figure in the results is about 130, cited from 1981. The Wik Me'anh people's territory was roughly 1,600 km$^2$ inland from Cape Keerweer. I found no school programme and no teaching materials.

\section{Primary Educational Resources}
No primer, textbook or school programme specific to Wik-Me'anha turned up in my searches.

\section{Linguistic and Reference Materials}
\begin{itemize}
\item Wikipedia article on the Wik-Me'nh language (a stub): \url{https://en.wikipedia.org/wiki/Wik-Me%27nh_language}
\item Wikipedia article on the Wik Me'anh people: \url{https://en.wikipedia.org/wiki/Wik_Me%27anh}
\item Wikipedia article on the Wik peoples, for regional context: \url{https://en.wikipedia.org/wiki/Wik_peoples}
\end{itemize}
Peter Sutton's 1978 PhD thesis on the Wik at Cape Keerweer is cited as the main scholarly source. I saw it only as a citation, with no URL, so it is not in the JSON.

\section{Religious and Community Resources}
None specific to Wik-Me'anha found.

\section{Digital and Audio Resources}
I found no audio, video or corpus specific to Wik-Me'anha.

\section{Curriculum Development Context}
Little is documented and the speaker figure is more than 40 years old. A curriculum would need new documentation work with Wik elders, probably at Aurukun, before any language-medium teaching is possible.
\end{document}